\appendix
\section{Reproducibility and artifact navigation}
\label{app:reproduction}
\begingroup\raggedright

The project retains executable definitions, experiment-specific manifests,
raw responses or measurements, classifier outputs, audits and rendered
reports. The canonical local workspace is
\path{/Users/yi/Documents/code/music}. Paths beginning with
\path{artifacts/} are relative to that root. Large audio and intermediate
products can reside on the remote NFS data volume rather than the local
report mirror. A local report's existence does not establish that its full
remote audio dependency graph is locally available.

\subsection{Data identities, views and storage}
The September 5 registry summary is
\path{artifacts/source_diversity_expansion_20260905/results/inventory/corpus_inventory.json},
with readable companion \path{corpus_inventory.md} in the same directory.
It records a particular frozen population; it is not a live total formed by
adding later overlapping cohorts. The local inventory snapshot is
\path{artifacts/audio_phenomena_expansion_20260907/audit/closeout_local_inventory_20260908_v1/files.jsonl}.
It lists files, not independent songs or a complete remote mirror.

Remote expansion code and logs use
\path{/mnt/nfs-code/users/yi/audio_phenomena_expansion_20260907};
its primary audio/measurement products use
\path{/mnt/nfs-data/users/yi/audio_phenomena_expansion_20260907}.
These paths are remote filesystem references, not links to local files.
Reproduction requires access to each manifest-bound source and compliance
with that source's access conditions; this thesis does not redistribute
restricted recordings or turn file possession into redistribution permission.

\subsection{Completed classification tables}
The following entry points retain full lattices rather than only the
selected comparisons displayed in the main text. Within
\path{artifacts/audio_phenomena_expansion_20260907/}:
\begin{description}
\item[Exact10.] \path{results/presentation_sdrfh_10s_v2/all_31x5_numeric_results.csv}
contains the 31-subset, five-cap numeric presentation.
\path{results/presentation_sdrfh_10s_v2/RESULTS_SDRFH_10S_EN.md}
provides the interpretation; \path{UNCERTAINTY_SDRFH_10S_RESULTS_EN.md}
documents the separate conditional and source-composition uncertainty analyses.
\item[Equal60.] \path{results/presentation_equal60_v6_v1/primary_macro_overview_255.csv}
contains all 255 subsets. \path{primary_matched_sc_deltas.csv} in that
directory retains matched SC increments, including negative values.
\path{primary_source_endpoints_all_arms.csv} supplies source endpoints,
and \path{training_row_group_ranges.csv} describes actual training support.
Files prefixed \path{diagnostic_all_cap_} preserve the separately labelled
diagnostic modes. The directory's \path{COMMIT.json} binds the presentation
products; it does not replace the model-level numerical audit.
\end{description}

The earlier four-family figures are
\path{artifacts/source_diversity_expansion_20260905/results/verified_presentation/10s_15x5_auc_ba.png}
and the corresponding \path{30s_15x5_auc_ba.png}. Their different cohorts
and feature/evaluation versions must be preserved in any cross-table comparison.

\subsection{Measurement evidence and scope of audits}
The external BC report is
\path{artifacts/audio_phenomena_expansion_20260907/results/bc_reserved_verified_report_v2/REPORT_EN.md},
with its retained \path{REPORT.tex} and rendered
\path{output/pdf/bc_reserved_verified_closeout_en.pdf} under the expansion
root. Producer completion, independent numerical reconstruction and the
scientific admission decision are separate evidence objects. Earlier failed
audit records remain part of the history; the corrected storage-layout
replay does not erase the original failure or change scientific tolerances.

The Native30 FHSC integrity report is
\path{artifacts/audio_phenomena_expansion_20260907/audit/native30_fhsc_committed_independent_audit_v1/REPORT_EN.md}.
Its file, schema, lineage and missingness checks must not be described as
independent re-extraction of every waveform descriptor. Conversely, a unit
test of a synthetic fixture is not evidence that the actual cohort completed.
Each claim in the thesis uses the audit level applicable to that claim.

\subsection{Execution and interpretation invariants}
Reproduction follows the dependency graph, not directory modification times:
verify source/view identities; verify frozen inference evidence; verify
per-recording measurements; assemble matched feature rows; verify the
fixed split/cap schedule; fit under a separate frozen evaluation contract;
replay predictions and endpoint aggregation; then render tables. An
experiment \path{COMMIT.json} denotes its immutable product manifest,
not necessarily a Git commit. A metadata schedule does not authorize fitting
and is not a count of completed models.

The CPU numerical path retains its frozen runtime and code dependencies.
Neural inference in this project was restricted to RTX 5090 hardware;
deterministic extraction and linear classification do not require a neural
GPU call. Failed attempts, scientific nulls and processing errors are kept
distinct. Resumption verifies existing products and does not replace failed
values by zero, silently change a split, or select a more favourable window.

Only training rows determine imputation, scaling and fitted coefficients.
Source/group weights, ridge penalty 10 and threshold 0.5 remain fixed.
Reproduction must retain all modes and omissions, preserve transitive group
purges, and avoid pooling raw scores from different fitted models into an
unqualified global AUC. Ordinary-group results, held-Human-source results
and held-generator results answer different questions.

\subsection{Document build and completed evaluation bindings}
The working master is
\path{interpretable_audio_thesis_working_v7_20260912_en.tex}.
From its containing directory, the retained build command is
\begin{verbatim}
latexmk -pdf -interaction=nonstopmode -halt-on-error \
  interpretable_audio_thesis_working_v7_20260912_en.tex
\end{verbatim}
Successful compilation is necessary but not sufficient: final acceptance
also requires resolved citations, traceable tables and visual inspection of
every page. The seven-family evaluation COMMIT SHA-256 is
\path{d1c16af22fa13b97fb2c9e69e6debc41b71c1d432cec86d1da035c06f3dd5efb}.
The report COMMIT SHA-256 is
\path{fecd7bef116dd7bfa0c3cfa8307c508be4fcb58acffbdc447011c3606c7ddb46}.
The completed BC, expanded YuE2 and Native60 transfer sections now provide
their terminal report/evaluation bindings and local result paths. The first two
report commits are respectively \texttt{ad5223a5f1f0fb66\ldots} and
\texttt{b7934f2c649e5d0f\ldots}; complete hashes are printed in those sections.
The Native60 prediction COMMIT is \texttt{34782ccfa467d530\ldots}.
Publication reconciliation, including redistribution restrictions for source
audio, is separate from experimental completion and remains in progress.
No prospective fit count is substituted for completion evidence.\par
\endgroup
